\sectionguard
\section{Reading Forward: The Rules of Typology}

What does the presence of God among men look like after Christ? The
Church's answer, so far as it concerns the ark, is her long habit of
reading the Virgin Mary as the ark of the new covenant. But before a
single parallel is laid on the table, the reader is owed a confession.
Resemblance is cheap. Set any two solemn processions side by side, in any
two literatures, and they will share a vocabulary: a going up, a shouting,
a carried holy thing, a crowd. Two texts written in the same Greek and
steeped in the same Scriptures will echo each other whether or not anyone
intended it. Whole libraries of nineteenth-century comparative religion
were built on parallels no weaker than some that will appear in the next
chapter, and those libraries are now cautionary tales. A reader who has
watched parallel-hunting done badly is right to keep his hand on his
wallet. Seeing a likeness is not yet knowing a meaning --- so it must be
said, once and briefly, what a likeness can prove.

\subsection{What Typology Is}

Typology is a claim about the authorship of history before it is a claim
about the art of any writer. In a well-made novel, the thing set quietly
on the mantel in the second chapter is aimed at the last chapter. No
analysis of the second chapter alone will explain it; only the author's
one intention, holding both chapters at once, explains it. Typology claims
that the two Testaments are such a book --- and that its prose is not
merely words but events, because its Author disposes events as freely as a
novelist disposes sentences.

This is not a private theory of the devout imagination. It is the Church's
stated rule of reading. The Second Vatican Council's constitution on
divine revelation puts the rule in two halves.\endnote{Second Vatican Council,
dogmatic constitution \work{Dei Verbum} (18 November 1965), \S 12, given
here in paraphrase, not quotation: the interpreter must carefully search out
the meaning the sacred writers really intended, attending to literary forms
and to the conditions of their time and culture; and, since sacred Scripture
must be read and interpreted in the same Spirit in which it was written,
must attend no less carefully to the content and unity of the whole of
Scripture, taking into account the living Tradition of the whole Church and
the analogy of faith.} The interpreter must first search out what the
human authors really intended, with every instrument that serves that
search; nothing that follows is excused from that duty. But because
Scripture is inspired, the same paragraph requires that it be read in the
same Spirit in which it was written, and names three tests of such
reading: the content and unity of the whole of Scripture; the living
Tradition of the whole Church; and the analogy of faith. The Catechism
receives the rule and gives the practice its name: from apostolic times
the Church has discerned, in the works of God under the old covenant,
prefigurations of what he accomplished in the fullness of time in his
incarnate Son.\endnote{\work{Catechism of the Catholic Church},
\S\S 128--130, in paraphrase; with \S\S 112--114, which restate \work{Dei
Verbum} 12's three criteria. \S 129 adds the two guards noted in the text:
Christians read the Old Testament in the light of Christ crucified and
risen, and such reading must not make the Old Testament's own standing as
revelation forgotten; the New Testament, equally, must be read in the light
of the Old.} Two words deserve their full weight. Typology
\emph{discerns}: the figure is found, not manufactured, and a discernment
can be mistaken, which is why evidence will matter at every step. And it
discerns in God's \emph{works}: deeds, persons, objects --- a chest of
wood as much as a word of prophecy --- because the Author of this book
writes with events.

It follows that typology presupposes the reality of what it reads; a
figure built on nothing signifies nothing. That is why the box itself came
first --- construction, journeys, catastrophe, disappearance --- and why
its questions of fact were settled as fact is settled. The canon models
the resulting reticence and the resulting confidence at once. The epistle
to the Hebrews inventories the ark's furniture --- the golden pot, the
rod, the tables --- and then closes the veil: ``of which it is not needful
to speak now particularly'' (Hebrews 9:5,
Douay-Rheims).\endnote{Hebrews 9:1--5,
Douay-Rheims (Challoner), read at \url{https://drbo.org/chapter/65009.htm};
v.~5 quoted.} The New Testament's own author signals that the furniture
means more than the letter has yet said, and declines to say it hastily.
There is no better warrant for speaking carefully.

Three confusions must be shut out at the gate. Typology is not literal
exegesis: the human author of 2 Kings 6 wrote of David's procession, and
no reading offered here will overwrite his letter. It is not a claim about
a writer's psychology: whether Luke consciously wove the ark narrative
into his first chapter is a real philological question, to be weighed with
the philologists' own tools; but the typological claim rests on the
intention of Another, so a demonstrated Lucan design would strengthen the
attestation, not create the meaning. And it is not decoration: a figure,
if true, teaches, and if false, deceives; in neither case is it ornament.

\subsection{The Two-Axis Rule}

Every disputed reading in what follows is handled by one named
distinction: \textbf{attestation is not reception, and the absence of
either is not falsity.} Attestation is the historical trail of a
reading --- who read the figure this way, when, in what words, under whose
name; a question of fact, settled by loci, dates, and manuscripts.
Reception is the Church's own adoption of a reading, in her liturgy, her
catechisms, her solemn teaching; equally a question of fact, settled by
dated public acts. The two run on separate axes. A reading may be
anciently attested yet never received into worship, or received at the
highest levels though no Father ever drew it; and a reading poor on both
axes may still be true, for late articulation is how development of
doctrine ordinarily looks from the outside. Accordingly, every reading
printed below wears its state on its sleeve: philologically anchored;
anciently attested, with its witnesses named and their names checked;
liturgically or magisterially received, with the act and its date; or a
labeled synthesis of this study's own.

One more instrument, because the next chapter leans on it from its first
page. The philological case there is cumulative: not one echo but a
cluster, gathered in a few verses. Each strand must be weighed alone, and
a strand that fails alone will be printed as failing; the negatives are
findings, not embarrassments. But the bundle must also be weighed as a
bundle, for the question a cluster raises is not whether any single echo
compels but how likely it is that so many should meet by accident in so
small a space. A rope is not refuted by showing that no one thread would
hold the load; a heap of threads is not yet a rope.

Two ditches lie beside this road, and travelers are in both. Consider
first the minimalist, and grant his conditional in full: if an allusion is
real only when it can be demonstrated --- shared words rare enough to
signify, in matching order, beyond coincidence --- then several of the
echoes in the next chapter fail the test. The critical school that says so
will be quoted there in its own words, at its own loci, with its premises
stated fairly.\endnote{The school --- Brown, Fitzmyer, Schürmann,
Bovon, with Isaacs's grounds --- is quoted verbatim and cited at exact loci
in the Visitation chapter below and in its notes; its shared premise, a
strict verbal-echo criterion for allusion, is stated there in the body.}
That work is honest and this argument needs it. The ditch lies one step
further on: the inference from ``the allusion is unproven as the
evangelist's intention'' to ``the reading is false.'' That step silently
assumes that the only meaning in Scripture is the human author's
demonstrable meaning --- the very question at issue. Philology can weigh
words; it cannot weigh the intentions of God. The minimalist who thinks it
can has not lost an argument; he has collapsed two axes into one.

Now the enthusiast, and grant his conditional too: if the Church's reading
is true, one would expect the ancients to have seen something --- and, in
substance, they did, earlier and more widely than the skeptic suspects, as
the Fathers will show in their place. But the enthusiast is not content
with the pedigree the record supplies; he improves it. Quotations migrate
under greater names. A sermon by no one becomes a sermon by a Doctor of
the Church. A translator's loose clause becomes a third-century witness.
An unchecked reading passes from book to book because it ought to be
true.\endnote{Each habit is documented at its place below: a sermon joining
David's dance before the ark to the Virgin, transmitted under two great
names, its one author unknown, and a celebrated ``earliest Marian ark''
text that proves an artifact of a nineteenth-century translation (the
chapter on the Fathers); and an ancient variant long repeated in apologetic
literature without an edition check --- which, as it happens, survived the
check when this study's research finally made it (the Visitation chapter).}
Notice what each embellishment concedes: that only an ancient pedigree
could make the reading true --- the minimalist's own axiom, swallowed
whole, with forgery added to make up the arrears. The Church does not need
the forgery and is not honored by it. Both ditches are one error: both
make attestation the sole measure of truth.

Why, finally, should reception count as evidence at all? Because the
council names the living Tradition of the whole Church among the criteria
of Spirit-led reading, and nowhere is that Tradition more public, more
deliberate, or more precisely datable than in worship; the manuals count
the liturgy among the \term{loci theologici} for exactly this reason.
When the Church appoints a lesson and a psalm to a feast, she is not
decorating a ceremony. She is performing an act of interpretation before
her whole people, and signing it with a date. A lectionary is a commentary
whose author is the reading Church, and its acts can be catalogued like
any other witnesses --- as, in their place, they will be.

That is all the method there is. The bar lifts, and the road beyond it
climbs --- as it happens, into the hill country of Juda.
